% ============================================================
% TUTORIAL: Fase 2 --- LLM Integration con SQS + Async Pattern
% Proyecto: Sentinel AcademIA
% Compilar: tectonic tutorial-fase2-llm.tex
% ============================================================

\documentclass[11pt,a4paper,oneside]{article}

\usepackage[utf8]{inputenc}
\usepackage[T1]{fontenc}
\usepackage[spanish,es-nodecimaldot]{babel}
\usepackage{lmodern}
\usepackage[margin=2.5cm]{geometry}
\usepackage{parskip}
\usepackage{microtype}

\usepackage{xcolor}
\usepackage{colortbl}
\usepackage{array}
\usepackage{booktabs}
\usepackage{amssymb}
\usepackage{pifont}
\usepackage{tcolorbox}
\tcbuselibrary{breakable,skins}

\usepackage{listings}

\lstdefinelanguage{yaml}{
  basicstyle=\ttfamily\small,
  keywordstyle=\color{blue!70!black}\bfseries,
  stringstyle=\color{red!70!black},
  commentstyle=\color{green!50!black}\itshape,
  morekeywords={AWSTemplateFormatVersion,Transform,Description,Parameters,Globals,Resources,Outputs,Type,Properties,Environment,Variables,Events},
}

\lstdefinelanguage{bash}{
  basicstyle=\ttfamily\small,
  keywordstyle=\color{blue!70!black}\bfseries,
  stringstyle=\color{red!70!black},
  commentstyle=\color{green!50!black}\itshape,
  morekeywords={aws,sam,cd,ls,mkdir,rm,cp,mv,export,unset,echo,cat,grep,curl,sleep,jq},
}

\lstdefinelanguage{json}{
  basicstyle=\ttfamily\small,
  showstringspaces=false,
  breaklines=true,
  literate=
    *{\{}{{{\{}}}{1}
     {\}}{{{\}}}}{1}
     {[}{{{[}}}{1}
     {]}{{{]}}}{1}
     {:}{{{:}}}{1}
     {,}{{{,}}}{1},
  string=[s]{"}{"},
  stringstyle=\color{red!70!black},
  morestring=[b]",
}

\lstset{
  basicstyle=\ttfamily\small,
  numberstyle=\tiny\color{gray},
  numbers=left,
  numbersep=8pt,
  frame=single,
  framerule=0.4pt,
  rulecolor=\color{gray!50},
  backgroundcolor=\color{gray!5},
  breaklines=true,
  breakatwhitespace=true,
  showstringspaces=false,
  captionpos=b,
  tabsize=2,
  literate={á}{{\'a}}1 {é}{{\'e}}1 {í}{{\'i}}1 {ó}{{\'o}}1 {ú}{{\'u}}1
           {ñ}{{\~n}}1 {Á}{{\'A}}1 {É}{{\'E}}1 {Í}{{\'I}}1 {Ó}{{\'O}}1 {Ú}{{\'U}}1
           {¿}{{?`}}1 {¡}{{!`}}1
}

\usepackage[colorlinks=true,linkcolor=blue!60!black,urlcolor=blue!60!black,citecolor=blue!60!black]{hyperref}

\usepackage{fancyhdr}
\pagestyle{fancy}
\fancyhf{}
\fancyhead[L]{\small\textsc{Sentinel AcademIA}}
\fancyhead[R]{\small Tutorial Fase 2}
\fancyfoot[C]{\thepage}
\renewcommand{\headrulewidth}{0.4pt}

\definecolor{okgreen}{RGB}{40,140,60}
\definecolor{errred}{RGB}{200,50,50}
\definecolor{warnorange}{RGB}{220,130,30}
\definecolor{infoblue}{RGB}{30,90,180}

\newcommand{\ok}{\textcolor{okgreen}{\textbf{\checkmark}}}
\newcommand{\err}{\textcolor{errred}{\textbf{\ding{55}}}}
\newcommand{\warn}{\textcolor{warnorange}{\textbf{\textasciitilde}}}

\title{\Huge\textbf{Fase 2: LLM Integration} \\[0.3em]
       \Large SQS + Async Pattern + Mock vs Real LLM \\[0.5em]
       \normalsize Tutorial paso a paso}
\author{Proyecto Sentinel AcademIA \\ \small lFunknown}
\date{\today}

\begin{document}

\maketitle
\thispagestyle{empty}

\begin{tcolorbox}[colback=blue!5!white,colframe=infoblue,title={\textbf{¿Qué construimos en este tutorial?}}]
En la \textbf{Fase 1} ten\'iamos un CRUD b\'asico (POST/GET/LIST de quejas). En esta fase agregamos \textbf{an\'alisis con IA} de forma as\'incrona:

\begin{itemize}
  \item \texttt{POST /api/quejas} persiste la queja y \textbf{encola} un mensaje
  \item Una Lambda \textbf{consumer de SQS} procesa la cola en background
  \item Esa Lambda llama al LLM (Cohere) y actualiza el status a \texttt{ANALIZADA}
  \item \texttt{GET /api/quejas/\{id\}} devuelve la queja con su an\'alisis
\end{itemize}

\textbf{Decisiones clave}:
\begin{itemize}
  \item \textbf{Por qu\'e SQS}: desacopla el POST r\'apido (202) del an\'alisis lento (LLM)
  \item \textbf{Por qu\'e Mock LLM en Lambda}: el package \texttt{oci} pesa 229MB y excede el l\'imite de 250MB de Lambda
  \item \textbf{Por qu\'e Real LLM local}: para demos completas con Cohere real
\end{itemize}
\end{tcolorbox}

\begin{tcolorbox}[colback=warnorange!5!white,colframe=warnorange,title={\textbf{Trampa de la fase: el l\'imite de 250MB de Lambda}}]
Cuando agregamos \texttt{pip install oci} a \texttt{requirements.txt}, el zip subi\'o a \textbf{292MB} y Lambda rechaz\'o el deploy con: \\
\texttt{Unzipped size must be smaller than 262144000 bytes}. \\
\textbf{Soluci\'on}: \texttt{MockLLMClient} (rule-based, sin deps externas) en Lambda, \texttt{RealLLMClient} solo local.
\end{tcolorbox}

\tableofcontents
\newpage

% ============================================================
\section{Arquitectura del flujo as\'incrono}
% ============================================================

\subsection{Diagrama del flujo}

\begin{center}
\begin{tabular}{|c|c|c|c|c|c|}
\hline
\textbf{Cliente} & $\rightarrow$ & \textbf{API Gateway} & $\rightarrow$ & \textbf{create\_queja} & $\rightarrow$ \\
\hline
POST /api/quejas & & & & Lambda & \\
\hline
\end{tabular}
\end{center}

\begin{center}
\begin{tabular}{|c|c|c|c|c|}
\hline
\textbf{DynamoDB} & $\leftarrow$ & \textbf{create\_queja} & $\rightarrow$ & \textbf{SQS} \\
\hline
Tabla Quejas & & persiste queja & & encola mensaje \\
\hline
\end{tabular}
\end{center}

\begin{center}
\begin{tabular}{|c|c|c|c|c|c|}
\hline
\textbf{SQS} & $\rightarrow$ & \textbf{analyze\_queja} & $\rightarrow$ & \textbf{LLM} & $\rightarrow$ \\
\hline
dispara & & Lambda consumer & & Mock o Cohere & \\
\hline
\end{tabular}
\end{center}

\begin{center}
\begin{tabular}{|c|c|c|c|c|}
\hline
\textbf{LLM response} & $\rightarrow$ & \textbf{analyze\_queja} & $\rightarrow$ & \textbf{DynamoDB} \\
\hline
Analisis JSON & & update status + analysis & & guarda todo \\
\hline
\end{tabular}
\end{center}

\subsection{Status workflow (state machine)}

\begin{center}
\begin{tabular}{lcl}
PENDIENTE & $\rightarrow$ & (POST recibido, en DynamoDB) \\
PENDIENTE & $\rightarrow$ & PROCESANDO (SQS triggered, an\'alisis iniciando) \\
PROCESANDO & $\rightarrow$ & ANALIZADA (LLM termin\'o, analysis guardado) \\
PROCESANDO & $\rightarrow$ & ERROR (LLM fall\'o 3 veces, va a DLQ) \\
ANALIZADA & $\rightarrow$ & NOTIFICADA (autoridad notificada, futuro) \\
\end{tabular}
\end{center}

\textbf{Tiempos esperados}:
\begin{itemize}
  \item POST $\rightarrow$ 202: $\sim$100ms (solo persiste + encola)
  \item PENDIENTE $\rightarrow$ ANALIZADA: $\sim$2-3s (SQS dispatch + LLM call)
\end{itemize}

% ============================================================
\section{Componente 1: SQS Queue + Dead Letter Queue}
% ============================================================

\subsection{Qu\'e es SQS}

\textbf{Amazon SQS} = cola de mensajes managed. Los producers env\'ian, los consumers reciben.

\textbf{Para qu\'e sirve en nuestro caso}:
\begin{itemize}
  \item \textbf{Desacoplar} el POST r\'apido del an\'alisis lento
  \item \textbf{Resiliencia}: si la Lambda consumer cae, los mensajes esperan
  \item \textbf{Retry autom\'atico}: si falla, SQS reintenta hasta 3 veces
  \item \textbf{DLQ} (Dead Letter Queue): mensajes que fallan 3 veces van ac\'a para an\'alisis manual
\end{itemize}

\subsection{Configuraci\'on en SAM (template.yaml)}

\begin{lstlisting}[language=yaml,caption={infra/template.yaml --- SQS}]
# Dead Letter Queue (mensajes que fallan 3 veces)
AnalysisDeadLetterQueue:
  Type: AWS::SQS::Queue
  Properties:
    QueueName: !Sub "sentinel-analysis-dlq-${Stage}"
    MessageRetentionPeriod: 1209600  # 14 dias

# Cola principal
AnalysisQueue:
  Type: AWS::SQS::Queue
  Properties:
    QueueName: !Sub "sentinel-analysis-${Stage}"
    VisibilityTimeout: 180  # 3 min para procesar
    RedrivePolicy:
      deadLetterTargetArn: !GetAtt AnalysisDeadLetterQueue.Arn
      maxReceiveCount: 3  # 3 reintentos antes de DLQ
\end{lstlisting}

\textbf{Conceptos}:
\begin{itemize}
  \item \texttt{VisibilityTimeout}: tiempo que un mensaje est\'a "invisible" despu\'es de ser le\'ido. Si la Lambda no termina en 180s, el mensaje vuelve a la cola.
  \item \texttt{maxReceiveCount: 3}: despu\'es de 3 entregas fallidas, va al DLQ.
  \item \texttt{DLQ Retention: 14 d\'ias}: tiempo que se guarda el mensaje muerto para investigaci\'on.
\end{itemize}

\subsection{Event mapping (SQS $\rightarrow$ Lambda)}

\begin{lstlisting}[language=yaml,caption={infra/template.yaml --- Lambda consumer}]
AnalyzeQuejaFunction:
  Type: AWS::Serverless::Function
  Properties:
    Role: !Ref LambdaRoleArn
    FunctionName: !Sub "sentinel-analyze-queja-${Stage}"
    CodeUri: ../src
    Handler: handlers.analyze_queja.lambda_handler
    MemorySize: 1024  # mas memoria para el LLM
    Timeout: 120       # mas tiempo para Cohere
    Environment:
      Variables:
        DYNAMODB_TABLE: !Sub "sentinel-quejas-${Stage}"
        MOCK_LLM: "true"  # usa MockLLMClient (no oci package)
    Events:
      SQSEvent:
        Type: Sqs
        Properties:
          Queue: !GetAtt AnalysisQueue.Arn
          BatchSize: 5  # procesa hasta 5 mensajes juntos
    Policies:
      - SQSPollerPolicy:
          QueueName: !GetAtt AnalysisQueue.QueueName
\end{lstlisting}

% ============================================================
\section{Componente 2: Analisis schema (Pydantic)}
% ============================================================

\subsection{El schema alineado con OpenAPI}

\begin{lstlisting}[language=Python,caption={handlers/\_shared/schemas.py --- Analisis}]
from enum import StrEnum

class CriticidadEnum(StrEnum):
    BAJA = "BAJA"
    MEDIA = "MEDIA"
    ALTA = "ALTA"
    CRITICA = "CRITICA"

class SentimientoEnum(StrEnum):
    POSITIVO = "POSITIVO"
    NEUTRO = "NEUTRO"
    NEGATIVO = "NEGATIVO"
    NEGATIVO_FUERTE = "NEGATIVO_FUERTE"

class Analisis(BaseModel):
    categoria: CategoriaEnum
    subcategoria: str | None = None
    criticidad: CriticidadEnum
    criticidadJustificacion: str | None = None
    sentimiento: SentimientoEnum
    entidades: list[Entidad] = Field(default_factory=list)
    temasClave: list[str] = Field(default_factory=list)
    accionSugerida: str
    prioridad: int = Field(ge=1, le=10)
    requiereNotificacionInmediata: bool = False
    modeloUsado: str = "cohere.command-latest"
    tokensConsumidos: int = 0
    latenciaMs: int = 0
    confidence: float = Field(ge=0, le=1, default=0.0)
    generatedAt: str
\end{lstlisting}

\subsection{El mensaje SQS}

\begin{lstlisting}[language=Python,caption={handlers/\_shared/schemas.py --- SQS message}]
class SQSCreateQuejaMessage(BaseModel):
    """Cuerpo del mensaje que create_queja envia a SQS."""
    quejaId: str
    tenantId: str
    correlationId: str | None = None
\end{lstlisting}

\textbf{Por qu\'e m\'inimo}: la Lambda consumer har\'a un \texttt{get\_queja()} a DynamoDB para obtener toda la info. El mensaje solo lleva el ID.

% ============================================================
\section{Componente 3: create\_queja ahora env\'ia a SQS}
% ============================================================

\subsection{El nuevo flujo}

\begin{lstlisting}[language=Python,caption={handlers/create\_queja.py --- agregar SQS}]
import boto3
from handlers._shared.config import get_settings

def _send_to_sqs(queja_id: str, tenant_id: str, correlation_id: str) -> None:
    """Envia mensaje a SQS para analisis async."""
    settings = get_settings()
    queue_url = settings.analysis_queue_url
    if not queue_url:
        log.warning("ANALYSIS_QUEUE_URL not configured, skipping SQS send")
        return

    sqs = boto3.client("sqs")
    message_body = json.dumps({
        "quejaId": queja_id,
        "tenantId": tenant_id,
        "correlationId": correlation_id,
    })
    sqs.send_message(
        QueueUrl=queue_url,
        MessageBody=message_body,
        MessageAttributes={
            "tenantId": {"StringValue": tenant_id, "DataType": "String"},
        },
    )
\end{lstlisting}

\subsection{Integraci\'on en el handler}

\begin{lstlisting}[language=Python,caption={lambda\_handler --- paso 6}]
# 5. Persistir en DynamoDB
dynamo = get_dynamo_client()
dynamo.put_item(item=item, condition_expression="attribute_not_exists(pk)")

# 6. Enviar a SQS (analisis async)
try:
    _send_to_sqs(accepted.quejaId, tenant_id, correlation_id)
except Exception as e:
    # No fallar el request: la queja ya esta persistida
    log.error("Failed to send SQS message (queja is still saved)",
              extra={"error": str(e)})

# 7. Return 202 Accepted
return HttpResponse.accepted(accepted.model_dump(mode="json"))
\end{lstlisting}

\textbf{Patrones importantes}:
\begin{itemize}
  \item \ok \textbf{try/except en SQS send}: si falla, la queja sigue guardada (mejor DX)
  \item \ok \textbf{MessageAttributes}: permiten filtrar en consumer si es necesario
  \item \ok \textbf{MessageBody es JSON}: estructurado, validable con Pydantic
\end{itemize}

% ============================================================
\section{Componente 4: analyze\_queja Lambda (SQS consumer)}
% ============================================================

\subsection{El handler completo}

\begin{lstlisting}[language=Python,caption={handlers/analyze\_queja.py}]
from handlers._shared.config import get_settings
from handlers._shared.correlation_id import (
    clear_correlation_id, with_correlation_id
)
from handlers._shared.dynamo_client import (
    get_dynamo_client, now_iso, queja_pk, queja_sk
)
from handlers._shared.errors import error_handler
from handlers._shared.logger import get_logger
from handlers._shared.schemas import QuejaResponse, SQSCreateQuejaMessage
from services.llm_client import LLMError, get_llm_client

log = get_logger(__name__)


def _set_status(queja_id: str, tenant_id: str, status: str) -> None:
    """Update queja status in DynamoDB."""
    dynamo = get_dynamo_client()
    dynamo.update_item(
        pk=queja_pk(tenant_id, queja_id),
        sk=queja_sk(queja_id),
        updates={"status": status, "updatedAt": now_iso()},
    )


def _process_one(record: dict) -> None:
    """Process a single SQS record."""
    body = record.get("body", "{}")
    msg = SQSCreateQuejaMessage.model_validate_json(body)
    log.append_keys(queja_id=msg.quejaId, tenant_id=msg.tenantId)

    # 1. Fetch queja
    dynamo = get_dynamo_client()
    item = dynamo.get_queja(tenant_id=msg.tenantId, queja_id=msg.quejaId)
    queja = QuejaResponse.from_dynamo(item)

    # 2. Update to PROCESANDO
    _set_status(msg.quejaId, msg.tenantId, "PROCESANDO")

    # 3. Call LLM
    llm = get_llm_client()
    try:
        analysis = llm.analyze_queja(queja)
    except LLMError as e:
        _set_status(msg.quejaId, msg.tenantId, "ERROR")
        raise  # SQS retry

    # 4. Update to ANALIZADA + store analysis
    analysis_json = analysis.model_dump(mode="json")
    dynamo.update_item(
        pk=queja_pk(msg.tenantId, msg.quejaId),
        sk=queja_sk(msg.quejaId),
        updates={
            "status": "ANALIZADA",
            "updatedAt": now_iso(),
            "analysis": analysis_json,
        },
    )


@error_handler
def lambda_handler(event, context):
    """SQS event handler for the analysis queue."""
    with_correlation_id(event)
    try:
        records = event.get("Records", [])
        failed = []
        for record in records:
            try:
                _process_one(record)
            except Exception:
                log.exception("Failed to process record")
                failed.append({
                    "itemIdentifier": record.get("messageId"),
                    "message": "error"
                })

        # batchItemFailures: SQS reintentar\'a solo los fallidos
        if failed:
            return {"batchItemFailures": failed}

        return {"statusCode": 200, "body": "OK"}
    finally:
        clear_correlation_id()
\end{lstlisting}

\subsection{Patrones importantes}

\begin{itemize}
  \item \ok \textbf{batchItemFailures}: SQS solo reintenta los fallidos (no todos)
  \item \ok \textbf{Status workflow}: PENDIENTE $\rightarrow$ PROCESANDO $\rightarrow$ ANALIZADA/ERROR
  \item \ok \textbf{log.append\_keys}: a\~nade quejaId/tenantId a TODOS los logs
  \item \ok \textbf{finally + clear}: limpia ContextVar al final
\end{itemize}

% ============================================================
\section{Componente 5: LLM Client (Mock + Real)}
% ============================================================

\subsection{El problema del 250MB}

Cuando agregamos \texttt{oci} a requirements.txt:

\begin{lstlisting}[language=bash,caption={Problema}]
$ du -sh .aws-sam-new/AnalyzeQuejaFunction/
292M  .aws-sam-new/AnalyzeQuejaFunction/   <-- EXCEDE 250MB

$ sam deploy ...
Error: Unzipped size must be smaller than 262144000 bytes
\end{lstlisting}

\textbf{El package \texttt{oci} pesa 229MB} (sus deps: cryptography, pycryptodome, etc.).

\subsection{La soluci\'on: factory con toggle}

\begin{lstlisting}[language=Python,caption={services/llm\_client.py --- factory}]
import os

def get_llm_client():
    """Get LLM client based on MOCK_LLM env var."""
    global _client
    if _client is not None:
        return _client

    use_mock = os.environ.get("MOCK_LLM", "true").lower() != "false"
    if use_mock:
        log.info("Using MockLLMClient (rule-based)")
        _client = MockLLMClient()
    else:
        log.info("Using RealLLMClient (OCI Cohere)")
        _client = RealLLMClient()
    return _client
\end{lstlisting}

\textbf{En SAM template}:
\begin{lstlisting}[language=yaml]
Environment:
  Variables:
    MOCK_LLM: "true"   # Lambda = mock
\end{lstlisting}

\textbf{En tu m\'aquina local}:
\begin{lstlisting}[language=bash]
unset MOCK_LLM         # o MOCK_LLM=false
python mi_script.py   # usa RealLLMClient
\end{lstlisting}

% ============================================================
\section{MockLLMClient: rule-based}
% ============================================================

\subsection{Algoritmo de clasificaci\'on}

\begin{lstlisting}[language=Python,caption={services/llm\_client.py --- MockLLMClient}]
class MockLLMClient:
    KEYWORD_MAP = {
        "ACOSO": (["acoso", "ofensiv", "agresi", "violencia",
                   "amenaza", "hostig"], "CRITICA", 10),
        "SALUD": (["suicid", "depres", "ansiedad", "salud mental"],
                  "CRITICA", 9),
        "INFRAESTRUCTURA": (["aula", "edificio", "baño", "techo",
                             "inund", "filtr", "electric", "dañado"],
                            "MEDIA", 5),
        "ACADEMICA": (["profesor", "docente", "calificaci",
                       "examen", "clase"], "MEDIA", 6),
        "ADMINISTRATIVA": (["matrícula", "inscripci", "trámite",
                            "sistema", "pago", "becas"], "MEDIA", 5),
    }

    def analyze_queja(self, queja: QuejaResponse) -> Analisis:
        text = f"{queja.titulo} {queja.descripcion}".lower()

        # 1. Detectar categoria por keywords
        best_match = ("OTRA", [], "BAJA", 3)
        for cat, (keywords, crit, prio) in self.KEYWORD_MAP.items():
            score = sum(1 for kw in keywords if kw in text)
            if score > 0 and (best_match[0] == "OTRA" or score > best_match[3]):
                best_match = (cat, keywords, crit, prio)

        # 2. Detectar sentimiento
        if any(w in text for w in ["terrible", "horrible", "urgente", "peligro"]):
            sentimiento = "NEGATIVO_FUERTE"
        elif any(w in text for w ["mal", "problema", "falla", "error"]):
            sentimiento = "NEGATIVO"
        # ... etc

        # 3. Build Analisis
        return Analisis(
            categoria=CategoriaEnum(categoria_str),
            criticidad=CriticidadEnum(criticidad_str),
            sentimiento=SentimientoEnum(sentimiento),
            prioridad=prioridad,
            accionSugerida=accion,
            # ...
        )
\end{lstlisting}

\textbf{Trade-offs}:
\begin{itemize}
  \item \ok \textbf{Ventaja}: 0 deps externas, 0 latencia, deterministic
  \item \err \textbf{Desventaja}: solo detecta keywords obvios, no entiende contexto
  \item \ok \textbf{Para el demo}: suficiente; el flujo async se demuestra igual
\end{itemize}

% ============================================================
\section{RealLLMClient: OCI Cohere (solo local)}
% ============================================================

\subsection{Cuando usarlo}

\begin{itemize}
  \item Para hacer demos en tu m\'aquina con el Cohere real
  \item Para tests que requieren LLM real
  \item NO en Lambda (excede 250MB)
\end{itemize}

\subsection{C\'odigo}

\begin{lstlisting}[language=Python,caption={services/llm\_client.py --- RealLLMClient}]
class RealLLMClient:
    def __init__(self):
        import oci
        from oci.generative_ai_inference import GenerativeAiInferenceClient
        from oci.generative_ai_inference import models as inference_models

        settings = get_settings()
        self.config = oci.config.from_file(settings.oci_key_file, "DEFAULT")
        # override con env vars si existen
        if settings.oci_user_ocid:
            self.config["user"] = settings.oci_user_ocid
        # ...

        self._client = GenerativeAiInferenceClient(
            self.config,
            service_endpoint=f"https://inference.generativeai.{self.region}.oci.oraclecloud.com",
        )
        self._inference_models = inference_models

    def analyze_queja(self, queja):
        user_prompt = self._build_user_prompt(queja)
        chat_request = self._inference_models.CohereChatRequest(
            message=user_prompt,
            max_tokens=800,
            temperature=0.2,
            is_stream=False,
        )
        details = self._inference_models.ChatDetails(
            compartment_id=self.compartment_id,
            serving_mode=self._inference_models.OnDemandServingMode(
                model_id=self.model_ocid,
            ),
            chat_request=chat_request,
        )
        response = self._client.chat(details)
        text = response.data.chat_response.text
        return self._parse_to_analisis(text)
\end{lstlisting}

\subsection{El prompt a Cohere}

\begin{lstlisting}[language=Python,caption={System prompt}]
SYSTEM_PROMPT = """Eres un asistente especializado en analizar quejas universitarias.
Devuelve SOLO JSON v\u00e1lido con esta estructura:

{
  "categoria": "ACADEMICA" | "INFRAESTRUCTURA" | "ACOSO" |
               "ADMINISTRATIVA" | "SALUD" | "OTRA",
  "criticidad": "BAJA" | "MEDIA" | "ALTA" | "CRITICA",
  "sentimiento": "POSITIVO" | "NEUTRO" | "NEGATIVO" | "NEGATIVO_FUERTE",
  "prioridad": 1-10,
  "requiereNotificacionInmediata": true | false,
  "confidence": 0.0-1.0,
  // ... etc
}

NO devuelvas markdown. SOLO el JSON puro."""
\end{lstlisting}

% ============================================================
\section{El bug de DynamoDB: floats no soportados}
% ============================================================

\subsection{El error real}

Cuando hicimos el primer test E2E, los logs mostraron:

\begin{lstlisting}[language=json,caption={CloudWatch logs --- error real}]
{
  "errorType": "TypeError",
  "errorMessage": "Float types are not supported. Use Decimal types instead.",
  "stackTrace": [
    "File \"/var/task/boto3/dynamodb/types.py\", line 171, in _is_number",
    "raise TypeError("
  ]
}
\end{lstlisting}

\subsection{Por qu\'e pas\'o}

El \texttt{Analisis.confidence} es un \texttt{float} (Python). DynamoDB \textbf{no acepta floats nativos}, solo \texttt{Decimal}:

\begin{lstlisting}[language=python]
# MAL
analysis.confidence = 0.75  # -> Python float
dynamo.update_item({"analysis": analysis.dict()})  # -> ERROR

# BIEN
from decimal import Decimal
analysis.confidence = Decimal("0.75")  # -> DynamoDB N
\end{lstlisting}

\subsection{El fix: helper \texttt{\_to\_dynamo}}

\begin{lstlisting}[language=Python,caption={handlers/\_shared/dynamo\_client.py}]
from decimal import Decimal

def _to_dynamo(value):
    """Recursively convert floats to Decimal for DynamoDB."""
    if isinstance(value, float):
        return Decimal(str(value))  # str() evita precision issues
    if isinstance(value, dict):
        return {k: _to_dynamo(v) for k, v in value.items()}
    if isinstance(value, list):
        return [_to_dynamo(item) for item in value]
    return value


def put_item(self, item, condition_expression=None):
    kwargs = {"Item": _to_dynamo(item)}  # <-- APLICAR
    # ...


def update_item(self, pk, sk, updates):
    # ...
    response = self._table.update_item(
        Key={"pk": pk, "sk": sk},
        UpdateExpression=update_expr,
        ExpressionAttributeValues=_to_dynamo(expression_values),  # <-- APLICAR
        # ...
    )
\end{lstlisting}

\textbf{Patrones}:
\begin{itemize}
  \item \ok \texttt{Decimal(str(x))} en vez de \texttt{Decimal(x)} para evitar \texttt{0.1 + 0.2 = 0.30000000000000004}
  \item \ok Recursivo: maneja dicts y lists anidados
  \item \err No lo olvides: si a\~nades otro client, tambi\'en debe usar \texttt{\_to\_dynamo}
\end{itemize}

% ============================================================
\section{Los 2 smoke tests E2E que validan todo}
% ============================================================

\subsection{Test 1: Queja de ACOSO (palabra clave)}

\begin{lstlisting}[language=bash,caption={Test E2E: ACOSO}]
API="https://om3eiczjr9.execute-api.us-east-1.amazonaws.com/dev"

# Paso 1: POST
RESP=$(curl -s -X POST "$API/api/quejas" \
  -H "Content-Type: application/json" \
  -H "X-Tenant-ID: utpl" \
  -H "X-User-ID: estudiante-003" \
  -H "X-Correlation-ID: fase2-final-001" \
  -d '{"titulo":"Situacion de acoso en el salon",
       "descripcion":"Varios companeros han reportado comentarios
                      agresivos y ofensivos reiterados. Situacion urgente.",
       "categoriaDeclarada":"ACOSO",
       "anonima":false}')

# Output esperado (202 Accepted)
# {"quejaId":"e3148f6c-...","status":"PENDIENTE",
#  "correlationId":"fase2-final-001",
#  "createdAt":"2026-06-20T14:02:45.077465+00:00",
#  "estimatedAnalysisTime":30}

# Paso 2: Esperar 5 segundos (SQS dispatch)
sleep 5

# Paso 3: GET con analysis
curl -s "$API/api/quejas/$QUEJA_ID" -H "X-Tenant-ID: utpl"
\end{lstlisting}

\begin{lstlisting}[language=json,caption={Output del GET --- analysis detectada}]
{
  "quejaId": "e3148f6c-92ec-4df2-9669-3539aeeb2195",
  "status": "ANALIZADA",
  "analysis": {
    "categoria": "ACOSO",
    "subcategoria": "acoso",
    "criticidad": "CRITICA",
    "criticidadJustificacion":
      "Detectadas palabras clave de ACOSO. Prioridad 10/10.",
    "sentimiento": "NEGATIVO_FUERTE",
    "entidades": [
      {"tipo": "MENCION", "texto": "Varios"},
      {"tipo": "MENCION", "texto": "Situacion"}
    ],
    "temasClave": ["acoso", "agresivos", "ofensivos"],
    "accionSugerida":
      "Notificar a BIENESTAR y SEGURIDAD inmediatamente.
       Escalar a direccion.",
    "prioridad": 10,
    "requiereNotificacionInmediata": true,
    "modeloUsado": "mock-rule-based-v1",
    "tokensConsumidos": 13,
    "latenciaMs": 0,
    "confidence": 0.75,
    "generatedAt": "2026-06-20T14:02:48.156952+00:00"
  }
}
\end{lstlisting}

\subsection{Test 2: Queja de INFRAESTRUCTURA}

\begin{lstlisting}[language=bash,caption={Test E2E: INFRAESTRUCTURA}]
RESP=$(curl -s -X POST "$API/api/quejas" \
  -H "Content-Type: application/json" \
  -H "X-Tenant-ID: utpl" \
  -H "X-Correlation-ID: fase2-infra-001" \
  -d '{"titulo":"Techo del salon con filtraciones",
       "descripcion":"El aula 205 tiene el techo dañado y se
                      inunda cuando llueve. Hay charcos constantes.",
       "categoriaDeclarada":"INFRAESTRUCTURA",
       "anonima":false}')

sleep 5
# Output del GET:
# {
#   "categoria": "INFRAESTRUCTURA",
#   "criticidad": "MEDIA",
#   "prioridad": 5,
#   "requiereNotificacionInmediata": false,
#   "temasClave": ["aula"],
#   "accionSugerida": "Asignar a comite academico. Responder
#                       en 5 dias habiles.",
#   "modeloUsado": "mock-rule-based-v1",
#   "confidence": 0.75
# }
\end{lstlisting}

% ============================================================
\section{Cat\'alogo de errores y soluciones}
% ============================================================

\begin{tcolorbox}[colback=errred!5!white,colframe=errred,breakable,title={Error 1: Unzipped size > 250MB}]
\textbf{Cuando}: \texttt{sam deploy} con \texttt{oci} en requirements. \\
\textbf{Causa}: el package \texttt{oci} pesa 229MB solo. \\
\textbf{Soluci\'on}: usar MockLLMClient en Lambda; oci solo local.
\end{tcolorbox}

\begin{tcolorbox}[colback=errred!5!white,colframe=errred,breakable,title={Error 2: Float types not supported}]
\textbf{Cuando}: guardar Analisis con \texttt{confidence: 0.75}. \\
\textbf{Causa}: DynamoDB no acepta \texttt{float}, solo \texttt{Decimal}. \\
\textbf{Soluci\'on}: helper \texttt{\_to\_dynamo} que convierte recursivamente.
\end{tcolorbox}

\begin{tcolorbox}[colback=errred!5!white,colframe=errred,breakable,title={Error 3: SQS message sent but consumer no triggered}]
\textbf{Cuando}: el mensaje va a SQS pero la Lambda no se ejecuta. \\
\textbf{Causa}: permisos del LabRole (no \texttt{sqs:ReceiveMessage}). \\
\textbf{Soluci\'on}: agregar \texttt{SQSPollerPolicy} en el template.
\end{tcolorbox}

\begin{tcolorbox}[colback=errred!5!white,colframe=errred,breakable,title={Error 4: batchItemFailures no reportado}]
\textbf{Cuando}: SQS reintenta TODOS los mensajes, no solo los fallidos. \\
\textbf{Causa}: la Lambda retorna 200 en vez de \texttt{\{batchItemFailures: [...]\}}. \\
\textbf{Soluci\'on}: capturar excepciones per-record y reportar failures.
\end{tcolorbox}

\begin{tcolorbox}[colback=errred!5!white,colframe=errred,breakable,title={Error 5: analysis en null despu\'es de esperar}]
\textbf{Cuando}: GET devuelve \texttt{analysis: null} tras 30s. \\
\textbf{Causa}: la Lambda fall\'o silenciosamente; revisar CloudWatch logs. \\
\textbf{Soluci\'on}: \texttt{aws logs tail /aws/lambda/sentinel-analyze-queja-dev}.
\end{tcolorbox}

% ============================================================
\section{Verificaci\'on en AWS Console}
% ============================================================

\subsection{Recursos desplegados (despu\'es de Fase 2)}

\begin{lstlisting}[language=bash,caption={Listar todos los recursos del stack}]
aws cloudformation describe-stack-resources \
  --stack-name sentinel-academia-dev \
  --query 'StackResources[].{Type:ResourceType, Id:PhysicalResourceId, Name:LogicalResourceId}' \
  --output table
\end{lstlisting}

\begin{lstlisting}[language=json,caption={Output esperado (resumen)}]
Type                    | Name                       | Status
------------------------|----------------------------|--------
AWS::DynamoDB::Table    | QuejasTable                | CREATE_COMPLETE
AWS::SQS::Queue         | AnalysisQueue              | CREATE_COMPLETE
AWS::SQS::Queue         | AnalysisDeadLetterQueue    | CREATE_COMPLETE
AWS::S3::Bucket         | FilesBucket                | CREATE_COMPLETE
AWS::Lambda::Function   | HealthCheckFunction        | CREATE_COMPLETE
AWS::Lambda::Function   | CreateQuejaFunction        | CREATE_COMPLETE
AWS::Lambda::Function   | GetQuejaFunction          | CREATE_COMPLETE
AWS::Lambda::Function   | ListQuejasFunction         | CREATE_COMPLETE
AWS::Lambda::Function   | AnalyzeQuejaFunction       | CREATE_COMPLETE
AWS::ApiGateway::RestApi| SentinelApi                | CREATE_COMPLETE
\end{lstlisting}

\subsection{Ver logs del an\'alisis en tiempo real}

\begin{lstlisting}[language=bash,caption={Tail logs de analyze\_queja}]
aws logs tail /aws/lambda/sentinel-analyze-queja-dev --follow
\end{lstlisting}

Output esperado:

\begin{lstlisting}[language=json]
{"level":"INFO","location":"lambda_handler:103",
 "message":"Processing 1 SQS records",
 "queja_id":"e3148f6c-92ec-4df2-9669-3539aeeb2195",
 "tenant_id":"utpl"}

{"level":"INFO","location":"_process_one:63",
 "message":"Status updated to PROCESANDO"}

{"level":"INFO","location":"get_llm_client:295",
 "message":"Using MockLLMClient (rule-based)"}

{"level":"INFO","location":"_process_one:84",
 "message":"Analysis stored",
 "categoria":"ACOSO","criticidad":"CRITICA","prioridad":10}
\end{lstlisting}

\subsection{Ver mensajes en DLQ (si hay errores)}

\begin{lstlisting}[language=bash,caption={Inspeccionar DLQ}]
aws sqs get-queue-attributes \
  --queue-url https://sqs.us-east-1.amazonaws.com/227165337884/sentinel-analysis-dlq-dev \
  --attribute-names All
\end{lstlisting}

% ============================================================
\section{Lo que aprendimos (resumen ejecutivo)}
% ============================================================

\begin{tcolorbox}[colback=okgreen!5!white,colframe=okgreen,title={\textbf{Checklist Fase 2}}]
\begin{enumerate}
  \item \textbf{SQS para desacoplar}: POST r\'apido, an\'alisis lento
  \item \textbf{DLQ para resiliencia}: mensajes fallidos van a cola aparte
  \item \textbf{VisibilityTimeout}: 3 min para procesar; si no, reintenta
  \item \textbf{batchItemFailures}: SQS solo reintenta los fallidos
  \item \textbf{Status workflow}: PENDIENTE $\rightarrow$ PROCESANDO $\rightarrow$ ANALIZADA
  \item \textbf{Mock vs Real LLM}: Toggle con \texttt{MOCK\_LLM} env var
  \item \textbf{Lambda 250MB limit}: \texttt{oci} package NO cabe; usar mock
  \item \textbf{DynamoDB Decimal}: floats de Python NO funcionan; helper \texttt{\_to\_dynamo}
  \item \textbf{log.append\_keys}: a\~nade context (quejaId, tenantId) a todos los logs
  \item \textbf{System prompt LLM}: instrucciones claras + JSON schema
\end{enumerate}
\end{tcolorbox}

% ============================================================
\section{Ejercicios para practicar}
% ============================================================

\subsection{Ejercicio 1: A\~nadir m\'as keywords al MockLLM}

Edita \texttt{MockLLMClient.KEYWORD\_MAP} para reconocer 5 categor\'ias m\'as. Test\'ealo con quejas reales.

\subsection{Ejercicio 2: Implementar RealLLMClient en Lambda con Layer}

Crea un Lambda Layer con el package \texttt{oci} slim (solo lo necesario: \texttt{oci.generative\_ai\_inference}). Adj\'untalo a \texttt{AnalyzeQuejaFunction}. Cambia \texttt{MOCK\_LLM=false}.

\subsection{Ejercicio 3: A\~nadir m\'etricas CloudWatch}

Mide cu\'anto tarda el LLM (latency), cu\'antas quejas se procesan por minuto, cu\'antas fallan. Usa \texttt{aws-lambda-powertools.metrics}.

\subsection{Ejercicio 4: Endpoint \texttt{GET /api/quejas/\{id\}/analysis}}

Crea un nuevo handler que devuelva SOLO el analysis (no toda la queja). \'Util para el dashboard.

\subsection{Ejercicio 5: A\~nadir notificaci\'on por email (Fase 4 prep)}

Cuando el status pasa a \texttt{NOTIFICADA}, env\'ia email v\'ia SES al admin. Usa el campo \texttt{requiereNotificacionInmediata} para filtrar.

% ============================================================
\section{Siguiente paso: Fase 3}
% ============================================================

En la \textbf{Fase 3} desplegamos el \textbf{frontend Vue 3}:

\begin{itemize}
  \item Configurar variable de entorno con la API URL
  \item Build de producci\'on con Vite
  \item Deploy a \textbf{Vercel} (o S3 + CloudFront)
  \item Probar el flujo end-to-end: frontend $\rightarrow$ backend $\rightarrow$ SQS $\rightarrow$ LLM
  \item Capturar screenshots para el demo
\end{itemize}

\textbf{Tiempo estimado}: 1-2 horas.

\begin{center}
\begin{tabular}{ll}
\toprule
Fase & Estado \\
\midrule
0. Setup + OCI & \checkmark \\
1. Backend Core & \checkmark \\
2. LLM integration (este tutorial) & \checkmark \\
3. Frontend deploy & pr\'oximo \\
4. Archivos + Email & pendiente \\
5. RAG & opcional \\
6. Testing & pendiente \\
7. Deploy final + demo & pendiente \\
\bottomrule
\end{tabular}
\end{center}

\vspace{2em}
\begin{center}
\textit{Fin del tutorial. ¡A por la Fase 3!}
\end{center}

\end{document}
